\chapter*{Image credits}
\markboth{Image credits}{Image credits}
\addcontentsline{toc}{chapter}{Image credits}
Historical portraits are credited beside each image. The d'Alembert portrait photograph is by Didier Descouens, CC BY-SA 4.0 (\url{https://creativecommons.org/licenses/by-sa/4.0/}); the Erd\H os photograph is by Topsy Kretts, CC BY 3.0 (\url{https://creativecommons.org/licenses/by/3.0/}). The two Wikimedia Commons file records are \href{https://commons.wikimedia.org/wiki/File:Maurice_Quentin_de_La_Tour_-_Portrait_de_Jean_Le_Rond_d'Alembert.jpg}{the d'Alembert portrait} and \href{https://commons.wikimedia.org/wiki/File:Erdos_head_budapest_fall_1992.jpg}{the Erd\H os photograph}. These licences govern the credited images and are exceptions to this book's general copyright notice.

The new photographs accompanying the workshops are NASA's portrait of Eileen Collins, NASA's Hubble eXtreme Deep Field, and NASA's photograph of the Apollo 11 launch. The Collins portrait and the Hubble image are the copies distributed with scikit-image 0.22.0; that documentation records them as public-domain NASA images. The Collins photograph was converted from PNG to JPEG; the Hubble image is unchanged. The launch photograph is NASA image 6900556, the Saturn V leaving the pad on 16 July 1969. NASA's media guidelines say that NASA images are generally not subject to copyright in the United States, and that a textbook may use them without separate permission when NASA is named and the picture is not offered as an endorsement. The source records are \href{https://www.flickr.com/photos/nasacommons/16504233985/}{NASA's Eileen Collins portrait}, \href{https://esahubble.org/news/heic1214/}{the Hubble eXtreme Deep Field release} and \href{https://images.nasa.gov/details-6900556}{NASA image 6900556}. The pictures provide context; they imply no endorsement of this book.

The mathematical diagrams in the new workshops are original vector drawings. A diagram illustrates a claim; its accompanying argument supplies the proof.
